\scanpage{17}
\mspage{15}
c'est un autre foncteur fibre, mais \uline{gradué}, donc défini par un \uline{torseur} $P$ \deleted{$\unsure{G}$} sous $G$, et $F$ \unsure{une forme|en termes} de $F'$ \unsure{---} par torseur $P'$ sous $G'$. Supposons que sur \deleted{\unsure{quelque} \unsure{point}} une extension $K'$ de $K$, on puisse trouver un isom $F\sim F'$ compatible avec les filtrations, et induisant l'identité sur les gradués associés.

[ \uline{Ex} \inserted{$S=\operatorname{Spec}k$,} $k$ de car. nulle, $F=T_{\mathrm{DR}}$ avec filtration canonique ; ça marche aussi si $S$ est n'importe quel schéma \inserted{régulier} de car. nulle \ldots ]. Alors il y a un sous-torseur \inserted{canonique} $Q$ de $P$ \deleted{$\unsure{G}$}, de groupe $H$ le sous-groupe de $G$ des autom compatibles avec filtration et induisant l'identité sur le gradué, et de $=$ \deleted{\unsure{pour}} un $Q'$ dans $P'$, \inserted{groupe $H'$.} On notera que $H,H'$ sont unipotents, Ainsi, $F$ se déduit canoniquement du foncteur gradué $F'$ [graduation définie par $\mathrm G_m\xrightarrow[K]{\alpha}G'$] à l'aide d'un torseur $Q'$ sous le groupe \inserted{unipotent} $H(\alpha)$ de $G'$ défini par $\alpha$ comme on pense.

% slanted
\margin{expliciter la construction de $H(\alpha)$}

13) \uline{Relations avec les groupes de \deleted{\unsure{Galois}}\inserted{Mumford-Tate}.}

Les structures de Hodge polarisables forment une $\otimes$-catégorie semi-simple. Toute sous-\deleted{\unsure{$\otimes$}}catégorie de type fini est ``définie'' par un \inserted{groupe alg. $G$ sur $Q$} \textsuperscript{\unsure{b}}.

Soit $(V,V_{\unsure{0|Q}},(V^{pq}))$ une structure de Hodge.

\textsuperscript{\unsure{b}} avec \unsure{$i'_1,i'_2$|$i_1,i_2$} : $\mathrm G_{m\,\unsure{\mathbb C|\mathrm C}}\unsure{-|\longrightarrow}G_{\mathbb C}$ satisfaisant $i_1(\lambda)i_2(\lambda)=i(\lambda)$ central, $i_2(\lambda)=\overline{i_1(\lambda)}$.

On a \unsure{réciproquement}
% archive sheet number: 16

